\chapter{Smart contracts}\label{ch:contracts}

The contracts are the part that makes the guarantee real. They are small on purpose (about 460 lines of Solidity 0.8.30 in total) and have no admin override, upgrade path or arbitrary-call entry point.

\section{The contracts}

\begin{stable}{@{}P{3.3cm}L@{}}
\headrow \thead{Contract} & \thead{Role}\\ \midrule
\code{MockUSDC} & Test stablecoin with 6 decimals and a hard supply cap of 1,000,000. Only the faucet can mint.\\
\code{MockUSDCFaucet} & Gives 1,000 MockUSDC per claim, at most once a day per address.\\
\code{LiquidVault} & Mock vault with no lockup. Illustrative yield 300 bps a year (metadata only: principal stays 1:1).\\
\code{TimedVault} & Mock vault with a real 30-day lockup. Illustrative yield 600 bps a year.\\
\code{BoundraFactory} & Fixes the reviewed token and the two vaults, and creates one executor per owner (\code{createExecutor}, \code{executorOf}).\\
\code{BoundraExecutor} & One per owner. Holds the owner's funds, stores the mandate, and is the only thing the agent (the \emph{operator}) can call.\\
\end{stable}

Both vaults have a \emph{risk state} set by the risk admin: \code{NORMAL} (0) or \code{DEGRADED} (1), plus a global withdrawal capacity. A degraded vault refuses new deposits. This is a controlled demo fixture that stands in for a real protocol incident, not an oracle.

\section{The mandate}

A mandate is a fixed tuple of twelve fields. The owner calls \code{setMandate(tuple)} on their executor; the contract records \code{keccak256(abi.encode(tuple))} as the mandate hash and emits \code{MandateSet(version, mandateHash, operator)}.

\begin{btable}{@{}P{4.1cm}P{2.6cm}L@{}}
\headrow \thead{Field} & \thead{Type} & \thead{Meaning}\\ \midrule\endhead
\code{chainId} & uint256 & Must be 84532 (Base Sepolia).\\
\code{asset} & address & The reviewed MockUSDC.\\
\code{operator} & address & The agent allowed to act. Boundra's agent or the owner's own.\\
\code{beneficiary} & address & Who receives withdrawals. Must equal the owner.\\
\code{expiresAt} & uint64 & When the mandate stops (seconds).\\
\code{allowedVaults} & bitmask & 1 Liquid, 2 Timed, 3 both.\\
\code{allowedActions} & bitmask & 1 deposit, 2 withdraw, 3 both.\\
\code{maxLockupSeconds} & uint32 & Longest lockup the owner accepted.\\
\code{maxDepositPerAction} & uint256 & Cap per deposit.\\
\code{maxCumulativeDeposits} & uint256 & Total deposit budget for this mandate.\\
\code{maxWithdrawPerAction} & uint256 & Cap per withdrawal.\\
\code{allocationTarget} & uint256 & Amount the single allocation aims for.\\
\end{btable}

\section{How value moves}

\begin{figure}[htbp]
\centering
\begin{tikzpicture}
  \node[boxlime,text width=2.6cm,minimum height=1.6cm] (owner) at (0,0) {\textbf{Owner wallet}};
  \node[boxgreen,text width=3.0cm,minimum height=1.6cm] (exec) at (5.6,0) {\textbf{Executor}\\\scriptsize idle funds + mandate};
  \node[box,text width=2.6cm,minimum height=1.6cm] (vault) at (11.2,0) {\textbf{Vault}\\\scriptsize principal 1:1};
  \node[boxdark,text width=2.6cm] (agent) at (5.6,-3.0) {\textbf{Operator}\\\scriptsize the agent};

  \draw[arr] ($(owner.east)+(0,0.42)$) -- node[lbl,above=1pt]{\stepno{1} approve + \code{fund}} ($(exec.west)+(0,0.42)$);
  \draw[arr] ($(exec.east)+(0,0.42)$) -- node[lbl,above=1pt]{\stepno{3} deposit} ($(vault.west)+(0,0.42)$);
  \draw[arr] ($(vault.west)+(0,-0.42)$) -- node[lbl,below=1pt]{\stepno{4} withdraw} ($(exec.east)+(0,-0.42)$);
  \draw[arr] ($(exec.west)+(0,-0.42)$) -- node[lbl,below=1pt]{\stepno{5} always to the owner} ($(owner.east)+(0,-0.42)$);
  \draw[arr,draw=bForest] (agent) -- node[lbl,right=4pt,align=left]{\stepno{2} \code{executeDeposit} or \code{executeWithdraw}\\only typed parameters, no recipient} (exec);
\end{tikzpicture}
\caption{The operator can only ask the executor to deposit or withdraw. Funds leave the system only to the owner.}
\end{figure}

\section{Rules the executor enforces}

Every operator call has one of two fixed shapes, with no recipient argument and no arbitrary calldata:

\begin{lstlisting}
executeDeposit (actionId, version, nonce, deadline, target, amount)
executeWithdraw(actionId, version, nonce, deadline, target, amount)
\end{lstlisting}

On each call the executor checks:

\begin{itemize}
  \item the caller is the mandate's operator, the mandate is not paused, revoked or expired, and the \code{version} is current;
  \item the \code{nonce} equals \code{nextActionNonce} and the \code{actionId} was never used (neither ever resets, even across mandates; a reverted call consumes neither);
  \item the \code{deadline} is in the future (exclusive) and not after the mandate's expiry;
  \item the target vault and the action are allowed, and the amount is within the per-action cap, the remaining budget and the allocation target;
  \item a deposit does not go into a degraded vault, and the lockup vault is only used if the mandate accepted its 30-day lockup and expires more than 30 days plus a one-day buffer after the deposit;
  \item there is \emph{exactly one} deposit allocation per mandate: a cap-limited partial deposit completes it, and the leftover stays idle and recoverable by the owner;
  \item a withdrawal pays the owner directly and never refills the deposit budget.
\end{itemize}

Owner controls survive everything: \code{pause}, \code{resume} (only if not expired or revoked), \code{revoke}, \code{recoverIdle(amount)} and \code{ownerWithdraw(target, amount)}. Recovery still respects the vault's real lockup and capacity. A new mandate can replace the old one only when both vault positions are zero, so a fresh deposit budget always needs the owner's explicit authorization after settlement.

\begin{caution}[What the contract cannot protect]
Within the mandate, a compromised operator can still act at inconvenient times or use the remaining budget in otherwise allowed ways. Caps limit damage; they do not prove that an action suits the owner financially. Already executed transactions are not undone by pause or revoke.
\end{caution}

\section{Events used as evidence}

\begin{stable}{@{}P{5.4cm}L@{}}
\headrow \thead{Event} & \thead{Used for}\\ \midrule
\code{ExecutorCreated(owner, executor)} & Registering the owner's executor (the API re-reads \code{executorOf}).\\
\code{MandateSet(version, mandateHash, operator)} & Matching the confirmed mandate to the approved preview.\\
\code{ActionExecuted(actionId, version, nonce, action, target, amount)} & Proof that an agent action really happened. A broadcast hash alone is not proof.\\
\code{RiskStateChanged(level, capacity)} & The worker's risk monitor, only from the two pinned vault addresses.\\
\code{Deposited}, \code{Withdrawn} & Vault-side position changes.\\
\end{stable}

\section{The Base Sepolia deployment}

Person 1 deployed the contracts with \code{contracts/script/DeployBoundra.s.sol} at block 48,322,785. Before it was pinned in \code{deploy/manifests/base-sepolia.json}, every role was read back from the contracts themselves and every contract's deployed bytecode was compared with the repository's compiled code. Because anyone can deploy look-alike contracts, the API and worker trust only this manifest: they check the chain, the addresses and the keccak256 hash of each contract's runtime code before acting.

\begin{stable}{@{}P{2.9cm}L@{}}
\headrow \thead{Contract} & \thead{Address}\\ \midrule
Factory & \texttt{0xFbD91368BfFA9AF5C7034bCB7F53Cc1eC906b7f4}\\
MockUSDC & \texttt{0x13e7a4d02941a029Cf841f16E38eF377cb21e1Cf}\\
Faucet & \texttt{0xEbdc09B132066fF067E30FE077A58E818801FA2f}\\
LiquidVault & \texttt{0x7C0fC1a0551d5099d093E8D8b34a214758f7B27a}\\
TimedVault & \texttt{0x7ffD08b7b52E26023142060D3f5dBDe27757CD84}\\
Risk admin & \texttt{0xA05F8266B63D8074AcBA1c8f4e7355E7F515F14b}\\
\end{stable}
